\chapter{The Efficiency Cost of Taxation}\label{ch:b3:efficiency-cost-of-taxation}

A government that spends a third of output must raise a third of output, and every
way of raising it changes what people buy, how much they work and what they report.
This chapter measures that cost: the \omterm{def:b3:efficiency-cost-of-taxation:harberger}{excess burden} of a tax and the formula that
approximates it, the cost of the last dalen raised, the rule that spreads taxes so
as to minimise the burden, who bears a tax in the end, and how much taxable income
responds to the rate --- the number on which the rest depends.

\begin{openingchart}{Central-government revenue as a share of GDP, six economies,
1870--2020 (percent). Sweden's 1996 value, a one-year spike to 41.8~\% in the
source between 22.2~\% and 31.7~\%, is replaced by the mean of its neighbours.}{Jordà, Schularick and Taylor, Macrohistory Database (revenue,
gdp); our ratio}
\begin{tikzpicture}
\begin{axis}[omtimesten, width=0.72\linewidth, height=0.44\linewidth, xmin=1865,
  xmax=2025, ymin=0, ymax=50, ylabel={Percent of GDP}, legend pos=north west,
  legend style={font=\footnotesize, legend columns=2}]
  \addplot[omDef, thick, mark=none] table[x=year, y=GBR] {figdata/bachelor-3/17-efficiency-cost-of-taxation/out/revenue.csv};
  \addplot[black, thick, mark=none] table[x=year, y=USA] {figdata/bachelor-3/17-efficiency-cost-of-taxation/out/revenue.csv};
  \addplot[omThm, thick, mark=none] table[x=year, y=FRA] {figdata/bachelor-3/17-efficiency-cost-of-taxation/out/revenue.csv};
  \addplot[omProp, thick, mark=none] table[x=year, y=SWE] {figdata/bachelor-3/17-efficiency-cost-of-taxation/out/revenue.csv};
  \addplot[omThm, thick, dashed, mark=none] table[x=year, y=ITA] {figdata/bachelor-3/17-efficiency-cost-of-taxation/out/revenue.csv};
  \addplot[omDef, thick, dashed, mark=none] table[x=year, y=JPN] {figdata/bachelor-3/17-efficiency-cost-of-taxation/out/revenue.csv};
  \legend{United Kingdom, United States, France, Sweden, Italy, Japan}
\end{axis}
\end{tikzpicture}
\end{openingchart}

% figdata: revenue_points
In 1913 central governments took between 2~\% of GDP (the United States) and 14~\%
(Japan). By 1950 the British share was 38~\%, and by 2015 the Swedish one 30~\% and the
Italian 28~\%. In Britain and the United States the rise came in steps, at the two
world wars; in Sweden and Italy it rose in peacetime: the Swedish share from 15~\% in
1950 to 29~\% in 1985, the Italian one by 13 points in the ten years to 1983. After
1950, only Japan's share fell back below its 1913 level.
Local governments and social-insurance funds, which the series leaves out, raise as
much again in several of these economies.

\begin{recall}
The grades 10--12 volume showed that a tax drives a wedge between the price buyers
pay and the price sellers receive, that its burden falls on the less elastic side
whatever the statute says, and that it creates a deadweight loss. The Year 1 volume
measured that loss with elasticities, and defined compensated demand, the
equivalent and compensating variations, and Lagrange's method.
\end{recall}

\begin{notation}[Symbols of this chapter]\label{not:b3:efficiency-cost-of-taxation:symbols}
$\tau$ is an ad valorem tax rate; $\varepsilon$ a compensated price elasticity, in absolute value;
$e$ the \omterm{def:b3:efficiency-cost-of-taxation:eti}{elasticity of taxable income} with respect to the net-of-tax rate $1 - \tau$; $a$ the
Pareto parameter of the top of the income distribution; $\lambda$ a Lagrange multiplier.
\end{notation}

\section{The excess burden of a tax}\label{sec:b3:efficiency-cost-of-taxation:burden}

A tax costs taxpayers more than the revenue it raises: they lose the tax they pay
and the trades it stops. The second part, the deadweight loss of earlier volumes, is
the tax's \emph{\omterm{def:b3:efficiency-cost-of-taxation:harberger}{excess burden}}. It must be measured on compensated demand: the income
effect of a tax is a transfer to the government, not a waste, and only the
substitution away from the taxed good is lost.

\begin{definition}[Harberger formula]\label{def:b3:efficiency-cost-of-taxation:harberger}
\index{Harberger formula}\index{excess burden}
The \emph{Harberger formula} approximates the excess burden of a small ad valorem tax
$\tau$ on a good with expenditure $pq$, supplied at constant cost, by $\tfrac12\varepsilon\tau^2 pq$, with $\varepsilon$ the
compensated elasticity of demand.
\end{definition}

\begin{proposition}[Harberger's approximation]\label{prop:b3:efficiency-cost-of-taxation:harberger}
With constant marginal cost $p$ and compensated demand linear near $q$, the \omterm{def:b3:efficiency-cost-of-taxation:harberger}{excess
burden} of a tax $\tau p$ per unit is $\tfrac12\varepsilon\tau^2 pq$: it rises with the square of the rate.
\end{proposition}

\begin{proof}
The tax raises the price from $p$ to $(1 + \tau)p$. Compensated demand falls by $\Delta q = \varepsilon q
\times \Delta p/p = \varepsilon\tau q$. The units no longer bought were worth to buyers between $p$ and
$(1 + \tau)p$, and cost $p$ to produce: the loss is the triangle of base $\Delta q$ and height $\tau p$,
$\tfrac12 \times \varepsilon\tau q \times \tau p = \tfrac12\varepsilon\tau^2 pq$.
\end{proof}

\begin{omfigure}
\begin{tikzpicture}[econ, x=0.55cm, y=0.55cm]
  \draw[econ axes] (0,0) -- (11,0) node[right, font=\scriptsize] {$q$};
  \draw[econ axes] (0,0) -- (0,8) node[above, font=\scriptsize] {price};
  \draw[curve, omDef] (1,7.5) -- (10,1) node[right, font=\scriptsize] {compensated demand};
  \draw[curve] (0,3) -- (10.5,3) node[right, font=\scriptsize] {$p$};
  \draw[curve, dashed] (0,4.8) -- (10.5,4.8) node[right, font=\scriptsize] {$(1 + \tau)p$};
  \fill[omThm!35] (4.74,4.8) -- (7.23,3) -- (4.74,3) -- cycle;
  \fill[omDef!15] (0,3) rectangle (4.74,4.8);
  \node[font=\scriptsize] at (2.4,3.9) {revenue};
  \node[font=\scriptsize, anchor=west] (eb) at (7.6,4.6) {excess burden};
  \draw[->, thin] (eb.west) -- (5.55,3.6);
  \draw[guide] (4.74,0) -- (4.74,3); \draw[guide] (7.23,0) -- (7.23,3);
  \node[below, font=\scriptsize] at (4.74,0) {$q_\tau$};
  \node[below, font=\scriptsize] at (7.23,0) {$q$};
\end{tikzpicture}
\par\smallskip{\small A tax $\tau p$ per unit: the rectangle is revenue, the triangle is
the \omterm{def:b3:efficiency-cost-of-taxation:harberger}{excess burden}. Doubling the rate doubles the triangle's height and base, so the
burden is four times larger.}
\end{omfigure}

The square matters. Raising the same revenue by two taxes at rate $\tau$ costs half as
much as one tax at $2\tau$ on a base of the same size; a broad base at a low rate beats a
narrow base at a high one.

\section{The marginal cost of public funds}\label{sec:b3:efficiency-cost-of-taxation:mcpf}

\begin{definition}[Marginal cost of public funds]\label{def:b3:efficiency-cost-of-taxation:mcpf}
\index{marginal cost of public funds}\index{marginal excess burden}
The \emph{marginal cost of public funds} is the cost to taxpayers of raising one more
dalen of revenue by a given tax: one plus the \emph{marginal \omterm{def:b3:efficiency-cost-of-taxation:harberger}{excess burden}}, the
increase in \omterm{def:b3:efficiency-cost-of-taxation:harberger}{excess burden} per dalen of extra revenue.
\end{definition}

\begin{proposition}[The marginal burden rises with the rate]\label{prop:b3:efficiency-cost-of-taxation:meb}
With a single tax of the Harberger proposition, revenue is $R = \tau pq(1 - \varepsilon\tau)$, the \omterm{def:b3:efficiency-cost-of-taxation:mcpf}{marginal
excess burden} is $\varepsilon\tau/(1 - 2\varepsilon\tau)$ and the \omterm{def:b3:efficiency-cost-of-taxation:mcpf}{marginal cost of public funds} is $(1 - \varepsilon\tau)/(1 - 2\varepsilon\tau)$,
both rising in $\tau$.
\end{proposition}

\begin{proof}
Taxed quantity is $q(1 - \varepsilon\tau)$, so $R = \tau pq(1 - \varepsilon\tau)$ and $dR/d\tau = pq(1 - 2\varepsilon\tau)$. The \omterm{def:b3:efficiency-cost-of-taxation:harberger}{excess burden}
$\tfrac12\varepsilon\tau^2pq$ has derivative $\varepsilon\tau pq$, linear in $\tau$. The ratio is $\varepsilon\tau/(1 - 2\varepsilon\tau)$; adding one gives
$(1 - \varepsilon\tau)/(1 - 2\varepsilon\tau)$. The numerator of the first rises and its denominator falls as $\tau$ rises.
\end{proof}

With $\varepsilon = 0.5$ and $\tau = 0.2$, the \omterm{def:b3:efficiency-cost-of-taxation:mcpf}{marginal excess burden} is $0.1/0.8 = 0.125$: the last dalen
costs taxpayers 1.125 dalens. A public project is worth funding at the margin only
if it returns more than the marginal cost of the funds that pay for it --- not more
than one dalen per dalen spent.

\section{The inverse-elasticity rule}\label{sec:b3:efficiency-cost-of-taxation:ramsey}

\begin{definition}[Ramsey rule (taxation)]\label{def:b3:efficiency-cost-of-taxation:ramsey-rule}
\index{Ramsey rule (taxation)}\index{inverse-elasticity rule}
The \emph{Ramsey rule} of taxation sets the tax rates on goods with independent
demands so as to raise a given revenue at the least \omterm{def:b3:efficiency-cost-of-taxation:harberger}{excess burden}; with
independent demands it takes the form of the \emph{inverse-elasticity rule}: each
rate is inversely proportional to the good's compensated elasticity.
\end{definition}

\begin{proposition}[Inverse elasticities]\label{prop:b3:efficiency-cost-of-taxation:ramsey}
With goods $i$ of pre-tax spending $x_i$ and independent compensated demands of
elasticity $\varepsilon_i$, minimising $\sum_i \tfrac12\varepsilon_i\tau_i^2x_i$ subject to revenue $\sum_i \tau_ix_i(1 - \varepsilon_i\tau_i) = R$ gives
$\tau_i = \lambda/\big((1 + 2\lambda)\varepsilon_i\big)$: rates in inverse proportion to elasticities.
\end{proposition}

\begin{proof}
The Lagrangian is $\sum_i \tfrac12\varepsilon_i\tau_i^2x_i - \lambda\big(\sum_i \tau_ix_i(1 - \varepsilon_i\tau_i) - R\big)$. Its derivative in $\tau_i$ is
$\varepsilon_i\tau_ix_i - \lambda x_i(1 - 2\varepsilon_i\tau_i) = 0$, so $\varepsilon_i\tau_i(1 + 2\lambda) = \lambda$ and $\tau_i = \lambda/\big((1 + 2\lambda)\varepsilon_i\big)$. The
factor $\lambda/(1 + 2\lambda)$ is the same for every good, so $\tau_i\varepsilon_i$ is equal across goods.
\end{proof}

The rule says to tax most what people cannot do without. Inelastic goods are often
necessities, so the efficient tax system is regressive, and every actual system
trades the rule off against the distribution of the burden. It also explains why
uniform value-added taxes are defended despite the rule: elasticities are not known
good by good, and every exemption invites lobbying for the next.

\section{Who bears the tax}\label{sec:b3:efficiency-cost-of-taxation:incidence}

\begin{definition}[Payroll tax]\label{def:b3:efficiency-cost-of-taxation:payroll-tax}
\index{payroll tax}
A \emph{payroll tax} is a tax levied on wages, paid by the employer, the employee or
both, often to finance social insurance.
\end{definition}

\begin{proposition}[The statutory side does not matter]\label{prop:b3:efficiency-cost-of-taxation:payroll}
In a competitive labour market with labour-supply elasticity $\eta_S$ and labour-demand
elasticity $\eta_D$ (absolute value), a small \omterm{def:b3:efficiency-cost-of-taxation:payroll-tax}{payroll tax} lowers the net wage by the
share $\eta_D/(\eta_S + \eta_D)$ of the tax and raises the employer's cost by $\eta_S/(\eta_S + \eta_D)$, whichever
side the law makes pay.
\end{proposition}

\begin{proof}
Let $w$ be the wage the worker receives and $w(1 + t)$ the employer's cost. Equilibrium
requires $S(w) = D\big(w(1 + t)\big)$. Differentiating at $t = 0$ in log terms, $\eta_S\,d\ln w = -\eta_D(d\ln w + dt)$, so
$d\ln w = -\eta_D\,dt/(\eta_S + \eta_D)$ and the employer's cost rises by $dt + d\ln w = \eta_S\,dt/(\eta_S + \eta_D)$. A tax
levied on the worker instead enters as $S\big(w(1 - t)\big) = D(w)$ and gives the same net wage
and cost to first order.
\end{proof}

\begin{omfigure}
\begin{tikzpicture}[econ, x=0.55cm, y=0.55cm]
  \draw[econ axes] (0,0) -- (11,0) node[right, font=\scriptsize] {employment};
  \draw[econ axes] (0,0) -- (0,8) node[above, font=\scriptsize] {wage};
  \draw[curve, omDef] (1,7.5) -- (10,1.5) node[right, font=\scriptsize] {labour demand};
  \draw[curve, omThm] (1,2) -- (10,6.2) node[right, font=\scriptsize] {labour supply};
  \draw[guide] (0,5.1) -- (4.6,5.1) node[pos=0, left, font=\scriptsize] {employer's cost};
  \draw[guide] (0,3.7) -- (4.6,3.7) node[pos=0, left, font=\scriptsize] {net wage};
  \draw[very thick] (4.6,3.7) -- (4.6,5.1);
  \node[font=\scriptsize, anchor=east] at (4.5,4.4) {tax wedge};
  \fill (5.85,4.27) circle (1.6pt);
  \node[font=\scriptsize, anchor=north west] at (5.95,4.15) {no tax};
\end{tikzpicture}
\par\smallskip{\small A \omterm{def:b3:efficiency-cost-of-taxation:payroll-tax}{payroll tax} opens a wedge between the employer's cost and the
net wage. The wedge sits where it fits between the curves, wherever the law places
it; the steeper (less elastic) curve takes the larger share.}
\end{omfigure}

The same logic runs beyond one market. A value-added tax on all consumption falls on
whatever people consume out of --- wages, profits, past savings --- and is a tax on
the income spent; a tax on the profits of a firm whose capital can move abroad
falls partly on the workers whose productivity that capital raised. What the law
names is the collector, not the payer.

\begin{omchart}{Central-government revenue by type, five economies, 1972 and 2024
(percent of revenue).}{World Bank, World Development Indicators (GC.TAX.YPKG.RV.ZS,
GC.TAX.GSRV.RV.ZS, GC.REV.SOCL.ZS, GC.TAX.INTT.RV.ZS); other is the remainder}
\begin{tikzpicture}
\begin{axis}[omrecord, xbar stacked, bar width=7pt, width=0.66\linewidth,
  height=0.5\linewidth, xmin=0, xmax=100, ymin=-0.7, ymax=9.7, ytick={0,...,9},
  yticklabels={United States 1972, 2024, United Kingdom 1972, 2024, France 1972,
    2024, Germany 1972, 2024, Korea 1972, 2024},
  y dir=reverse, yticklabel style={font=\scriptsize}, xlabel={Percent of revenue},
  grid=none, xmajorgrids, legend style={at={(0.5,-0.2)}, anchor=north,
  legend columns=3, font=\footnotesize},
  cycle list={{fill=omDef!80, draw=white}, {fill=omThm!70, draw=white},
    {fill=omProp!70, draw=white}, {fill=omMeth!70, draw=white},
    {fill=black!25, draw=white}}]
\pgfplotsforeachungrouped \col in {income, goods, social, trade, other} {
  \addplot+[] table[y=pos, x=\col]
    {figdata/bachelor-3/17-efficiency-cost-of-taxation/out/composition.csv};
}
\legend{Income and profits, Goods and services, Social contributions, Trade, Other}
\end{axis}
\end{tikzpicture}
\end{omchart}

% figdata: composition
The mix changes slowly and differs widely. Social contributions --- \omterm{def:b3:efficiency-cost-of-taxation:payroll-tax}{payroll taxes},
mostly --- rose from 24 to 37~\% of United States federal revenue between 1972 and
2024, and from 1 to 30~\% in Korea, where trade taxes fell from 10~\% to 1~\%. France
and Germany draw more than a third of central revenue from social contributions;
the United States draws more than half from income and profits.

\section{The elasticity of taxable income}\label{sec:b3:efficiency-cost-of-taxation:eti}

\begin{definition}[Elasticity of taxable income]\label{def:b3:efficiency-cost-of-taxation:eti}
\index{elasticity of taxable income}
The \emph{elasticity of taxable income} is the percentage change in the income
people report to the tax authority for a one percent change in the net-of-tax rate
$1 - \tau$.
\end{definition}

Taxable income responds through every margin at once: hours and effort, but also
the timing of income, its form (wages, dividends, capital gains), deductions,
avoidance and evasion. If the tax law is fixed, all of these are costs to the
economy of raising the rate, and the \omterm{def:b3:efficiency-cost-of-taxation:eti}{elasticity of taxable income} measures them in
one number.

\begin{model}{The elasticity of taxable income}
\assumes Reported income responds to the net-of-tax rate with a stable elasticity,
and the income that disappears from the base is not taxed elsewhere.
\predicts The \omterm{def:b3:efficiency-cost-of-taxation:harberger}{excess burden} of an income tax includes all responses, avoidance as
well as work; revenue from the top bracket peaks at a finite rate.
\tested Tax reforms set the reported incomes of the groups whose rates change
against groups whose rates do not: the U.S. 1993 top-rate increase, a sequence of
Danish reforms with population-wide data since 1980, the United Kingdom's 50~\% rate
of 2010--2013. % ledger: b3ec:ssg.review, b3ec:ks.denmark
\breaks The elasticity is not a constant of nature. It is large where the tax base
has holes --- income that can be shifted into capital gains or into another year ---
and small where the base is broad and enforcement strong; income that moves to
another base is not a full loss.
\end{model}

\begin{measure}[How an elasticity of taxable income is estimated]
Two designs dominate. A difference-in-differences compares the change in reported
income of those whose rate changed with the change for a nearby group whose rate
did not, and must deal with trends in inequality that differ between the groups and
with mean reversion (a high income this year tends to fall next year anyway). At a
kink of the schedule, where the marginal rate jumps, people who respond pile up just
below the kink; the excess number of taxpayers there, against a smooth density,
gives the elasticity. In United States tax returns, bunching shows up at the first
kink of the earned-income credit, but only among the self-employed, and at the
threshold where income tax starts, and not at the other kinks. % ledger: b3ec:saez.bunching
\end{measure}

\section{Revenue against the rate}\label{sec:b3:efficiency-cost-of-taxation:laffer}

\begin{definition}[Revenue-maximising tax rate; Laffer curve]\label{def:b3:efficiency-cost-of-taxation:laffer}
\index{revenue-maximising tax rate}\index{Laffer curve}
The \emph{Laffer curve} is tax revenue as a function of the tax rate, zero at rates of
0 and 100~\%. The \emph{revenue-maximising tax rate} is the rate at its peak.
\end{definition}

\begin{proposition}[The revenue-maximising top rate]\label{prop:b3:efficiency-cost-of-taxation:topRate}
Let incomes above a threshold $\bar z$ have mean $z_m$, with $a = z_m/(z_m - \bar z)$ constant (a Pareto tail),
and let $z_m$ respond to $1 - \tau$ with elasticity $e$. Revenue from a marginal rate $\tau$ on income
above $\bar z$ is maximised at $\tau^* = 1/(1 + ae)$.
\end{proposition}

\begin{proof}
Per taxpayer in the bracket, revenue is $\tau(z_m - \bar z)$. A small rise $d\tau$ raises it mechanically
by $(z_m - \bar z)\,d\tau$ and lowers it by $\tau\,dz_m = -\tau ez_m\,d\tau/(1 - \tau)$. At the peak the two balance:
$(1 - \tau)(z_m - \bar z) = \tau ez_m$, that is, $(1 - \tau)/a = \tau e$, and $\tau^* = 1/(1 + ae)$.
\end{proof}

\begin{omchart}{Revenue from a top bracket against its rate, for a Pareto tail with
$a = 1.5$ and elasticities of taxable income of 0.25 and 0.5 (the peak with $e = 0.25$
is 1).}{Our computation (\code{make\_laffer.py})}
\begin{tikzpicture}
\begin{axis}[omrecord, width=0.56\linewidth, height=0.36\linewidth, xmin=0, xmax=100,
  ymin=0, ymax=1.1, xlabel={Top marginal rate (\%)}, ylabel={Revenue},
  legend pos=north west, legend style={font=\footnotesize}]
  \addplot[omDef, thick, mark=none] table[x=tau, y=e25] {figdata/bachelor-3/17-efficiency-cost-of-taxation/out/laffer.csv};
  \addplot[omThm, thick, dashed, mark=none] table[x=tau, y=e50] {figdata/bachelor-3/17-efficiency-cost-of-taxation/out/laffer.csv};
  \legend{$e = 0.25$, $e = 0.5$}
\end{axis}
\end{tikzpicture}
\end{omchart}

% figdata: laffer_peaks
With $a = 1.5$, the revenue-maximising rate is 73~\% if the elasticity is 0.25 and 57~\%
if it is 0.5. The formula needs two numbers: $a$ is measured from tax returns with
little dispute; $e$ is the hard one, and the whole argument turns on it.

\begin{record}[The United Kingdom's 50~\% rate, 2010--2013]
A 50~\% rate on incomes above 150{,}000 pounds took effect in April 2010, announced
in March 2009, and was cut to 45~\% from April 2013. Tax officials estimated that
about 18 billion pounds of income had been brought forward into 2009--10, before
the rate rose, and an underlying elasticity of 0.48; the official costing of the cut
used 0.45, which put the extra revenue of the 50~\% rate over a 45~\% rate at 0.1
billion pounds. Later work with fuller data found 0.31 for the first year and 0.83
for the second, and between 0.1 and 0.2 for taxpayers just above the threshold.
% ledger: b3ec:ifs.50p
The episode shows what moves the elasticity: a rise announced a year ahead invites
retiming, and the richest, whose income is most easily shifted, respond most.
\end{record}

\section{Taxes people do not notice}\label{sec:b3:efficiency-cost-of-taxation:salience}

\begin{definition}[Tax salience]\label{def:b3:efficiency-cost-of-taxation:salience}
\index{tax salience}
\emph{Tax salience} is the visibility of a tax to the people who pay it at the moment
they decide: a tax included in the posted price is salient; one added at the till
is less so.
\end{definition}

\begin{record}[Posted and unposted sales taxes]
In a grocery store experiment, posting price tags that included the sales tax cut
demand for the tagged goods by 8~\%, though the price paid at the till was the same.
Across states, increases in alcohol taxes that are included in posted prices
reduced consumption more than increases in taxes added at the register.
% ledger: b3ec:clk.salience
\end{record}

Salience breaks two results of this chapter. The statutory side matters: a tax that
buyers do not see when they choose falls more on them than the elasticities
predict. And a tax that changes nothing people do can still cost them, because they
spend as if their income were larger than it is and end up with the wrong bundle;
the \omterm{def:b3:efficiency-cost-of-taxation:harberger}{excess burden} is no longer the triangle under a demand curve that ignores the
tax.

\section{Exercises}\label{sec:b3:efficiency-cost-of-taxation:exercises}

\begin{exercise}[$\star$]\label{exo:b3:efficiency-cost-of-taxation:1}
In the opening chart, when did central-government revenue rise most in each
economy? Did it ever fall back to its pre-1913 level?
\end{exercise}

\begin{exercise}[$\star$]\label{exo:b3:efficiency-cost-of-taxation:2}
A tax of 20~\% falls on a good with compensated elasticity 0.5 and spending of 100
million dalens. Compute the \omterm{def:b3:efficiency-cost-of-taxation:harberger}{excess burden}, the revenue and their ratio.
\end{exercise}

\begin{exercise}[$\star$]\label{exo:b3:efficiency-cost-of-taxation:3}
With the numbers of exercise 2, compute the \omterm{def:b3:efficiency-cost-of-taxation:mcpf}{marginal excess burden} and the \omterm{def:b3:efficiency-cost-of-taxation:mcpf}{marginal
cost of public funds}. A project returns 1.1 dalens per dalen spent: should it be
funded by raising this tax?
\end{exercise}

\begin{exercise}[$\star$]\label{exo:b3:efficiency-cost-of-taxation:4}
What evidence would contradict the claim that a value-added tax is passed on in
full to consumer prices?
\end{exercise}

\begin{exercise}[$\star\star$]\label{exo:b3:efficiency-cost-of-taxation:5}
Two goods with independent demands have compensated elasticities 0.2 and 0.8 and
equal spending. The government must raise 10~\% of that spending. Find the
Ramsey rates, using the \omterm{def:b3:efficiency-cost-of-taxation:ramsey-rule}{inverse-elasticity rule} and the revenue constraint
(solve numerically), and compare the \omterm{def:b3:efficiency-cost-of-taxation:harberger}{excess burden} with a uniform tax raising the
same revenue.
\end{exercise}

\begin{exercise}[$\star\star$]\label{exo:b3:efficiency-cost-of-taxation:6}
Find the revenue-maximising top rate with $a = 1.5$ and $e = 0.25$; then with $a = 2$. Why
does a thinner tail lower the rate?
\end{exercise}

\begin{exercise}[$\star\star$]\label{exo:b3:efficiency-cost-of-taxation:7}
Labour supply has elasticity 0.2 and labour demand 0.6. A \omterm{def:b3:efficiency-cost-of-taxation:payroll-tax}{payroll tax} of 10~\% of
the wage is levied on employers. By how much does the net wage fall, and the
employer's cost rise? What changes if the tax is levied on workers?
\end{exercise}

\begin{exercise}[$\star\star$]\label{exo:b3:efficiency-cost-of-taxation:8}
With the WDI data, compute for the five economies of the composition chart the
change between 1990 and 2024 in the share of each type of revenue. Which type gained
most, and where?
\end{exercise}

\begin{exercise}[$\star\star$]\label{exo:b3:efficiency-cost-of-taxation:9}
What evidence would show that bunching at a kink reflects changes in what people
earn, not only in what they report?
\end{exercise}

\begin{exercise}[$\star\star\star$]\label{exo:b3:efficiency-cost-of-taxation:10}
Prove the \omterm{def:b3:efficiency-cost-of-taxation:ramsey-rule}{inverse-elasticity rule} for $n$ goods with independent demands, and show
that with dependent demands (substitutes) the rule must change. In which direction
for a good that substitutes for leisure?
\end{exercise}

\begin{exercise}[$\star\star\star$]\label{exo:b3:efficiency-cost-of-taxation:11}
At a kink, the marginal rate rises from 20~\% to 35~\% at an income of 40{,}000. The
density of taxpayers just around the kink is 50 per 1{,}000 dalens, and 120 more
taxpayers are found within the bunching window than a smooth density would put
there. Find the income response of the marginal buncher, and the elasticity it
implies.
\end{exercise}

\begin{exercise}[$\star\star\star$]\label{exo:b3:efficiency-cost-of-taxation:12}
A sales tax of 10~\% is added at the till. Buyers react to it as if it were 4~\%.
Supply is perfectly elastic; the compensated elasticity of demand is 1 and spending
is 100. Who bears the tax, and what is the \omterm{def:b3:efficiency-cost-of-taxation:harberger}{excess burden}?
What changes if the tax is included in posted prices?
\end{exercise}

\section{Problem: Brenland's Tax Reform}\label{sec:b3:efficiency-cost-of-taxation:problem}

\begin{problem}[{Weekend problem --- Brenland's tax reform}]\label{pb:b3:efficiency-cost-of-taxation:brenland}
Brenland's income tax of 2025 is zero up to 12{,}000 dalens, 20~\% from 12{,}000 to 40{,}000,
35~\% from 40{,}000 to 100{,}000 and 45~\% above. The household survey of 2{,}000 households
(\code{load("households")}) records gross income and tax. The upper bands are the
income above 40{,}000. Take an \omterm{def:b3:efficiency-cost-of-taxation:eti}{elasticity of taxable income} of 0.25.
% figdata: brenland_stats

\medskip\noindent\textbf{Part I --- the schedule.}
\begin{enumerate}
\item Compute the tax and the average rate of a household earning 25{,}000, and of one
earning 60{,}000.
\item What are their marginal rates?
\item How many of the 2{,}000 households have income above 40{,}000, and what is their mean
income?
\item What do the upper bands raise from them (the tax on income above 40{,}000), and
what share of the total income tax is it?
\item Compute the Pareto parameter $a = z_m/(z_m - 40{,}000)$ of these households.
\item Brenland's top band has four households. Why is the analysis done on the upper
bands together?
\end{enumerate}

\medskip\noindent\textbf{Part II --- the burden of the upper bands.}
\begin{enumerate}[resume]
\item Every upper rate rises by one point. Compute the mechanical revenue gain.
\item Compute the behavioural loss, $e\sum\tau_iz_i/(1 - \tau_i)$ per point, over the households
above 40{,}000.
\item Compute the \omterm{def:b3:efficiency-cost-of-taxation:mcpf}{marginal cost of public funds} of the upper bands.
\item Compare with the formula $1/\big(1 - \tau ea/(1 - \tau)\big)$ at $\tau = 0.35$.
\item Compute the revenue-maximising rate for these bands.
\item Redo question 9 with $e = 0.5$. What does the result mean?
\end{enumerate}

\medskip\noindent\textbf{Part III --- a reform.} Brenland cuts every upper rate by two
points and raises its value-added tax of 18~\% on household consumption to replace
the lost revenue. Treat consumption as fixed, and take the \omterm{def:b3:efficiency-cost-of-taxation:mcpf}{marginal cost of public
funds} of the value-added tax as 1.25.
\begin{enumerate}[resume]
\item Compute the static revenue loss of the cut.
\item Compute the revenue regained by the behavioural response (first order).
\item Compute the net loss to be replaced.
\item By how many points must the value-added tax rate rise?
\item Compute, to first order, the gain to Brenland's taxpayers from the swap.
\item Who gains and who loses from the reform?
\item Why might the value-added tax's \omterm{def:b3:efficiency-cost-of-taxation:mcpf}{marginal cost of public funds} be smaller than
the upper bands'?
\end{enumerate}

\medskip\noindent\textbf{Part IV --- a project at the margin.} A public project
returns 1.5 dalens of benefits per dalen spent.
\begin{enumerate}[resume]
\item Should it be funded by raising the upper rates?
\item Should it be funded by raising the value-added tax?
\item What \omterm{def:b3:efficiency-cost-of-taxation:eti}{elasticity of taxable income} would make the upper bands' \omterm{def:b3:efficiency-cost-of-taxation:mcpf}{marginal cost
of public funds} equal to 1.5?
\item What evidence on Brenland's taxpayers would you want before using $e = 0.25$?
\item How do the United Kingdom's estimates for the 50~\% rate bear on that choice?
\item Why does the answer depend on the tax base as well as on the rate?
\item State the result: the \omterm{def:b3:efficiency-cost-of-taxation:mcpf}{marginal cost of public funds} of Brenland's upper bands.
\end{enumerate}
\end{problem}
